% ═══════════════════════════════════════════════════════════════════════════
% SLIDES: Repaso Unidad 1 + 2
% ═══════════════════════════════════════════════════════════════════════════

\documentclass[aspectratio=169,11pt]{beamer}

\usepackage[utf8]{inputenc}
\usepackage[T1]{fontenc}
\usepackage[ngerman]{babel}

% Theme
\usetheme{default}
\usecolortheme{default}

% Farben
\definecolor{spanisch}{HTML}{c41e3a}
\definecolor{hellgrau}{HTML}{f5f5f5}

\setbeamercolor{structure}{fg=spanisch}
\setbeamercolor{frametitle}{fg=spanisch,bg=white}
\setbeamercolor{title}{fg=spanisch}
\setbeamercolor{block title}{fg=white,bg=spanisch}
\setbeamercolor{block body}{fg=black,bg=hellgrau}

% Keine Navigation
\setbeamertemplate{navigation symbols}{}
\setbeamertemplate{footline}[frame number]

% Pakete
\usepackage{tikz}
\usepackage{qrcode}

\title{Repaso: Unidad 1 + 2}
\subtitle{Prüfungsvorbereitung}
\author{Klasse 9 -- Spanisch}
\date{}

\begin{document}

% ═══════════════════════════════════════════════════════════════════════════
% TITELFOLIE
% ═══════════════════════════════════════════════════════════════════════════

\begin{frame}
\titlepage
\end{frame}

% ═══════════════════════════════════════════════════════════════════════════
% ZIEL
% ═══════════════════════════════════════════════════════════════════════════

\begin{frame}{Hoy vamos a...}

\begin{block}{Ziel der Stunde}
Wiederholung aller wichtigen Inhalte aus Unidad 1 + 2
\end{block}

\vspace{5mm}

\textbf{Was ihr wiederholt:}
\begin{itemize}
    \item Vokabeln: Restaurante, Casa, Muebles
    \item Verben: pedir, servir, conocer, parecer
    \item Pretérito perfecto (regulär + irregulär)
    \item Direktobjektpronomen (lo/la/los/las)
    \item Bewegungsverben mit Präpositionen
\end{itemize}

\end{frame}

% ═══════════════════════════════════════════════════════════════════════════
% ABLAUF
% ═══════════════════════════════════════════════════════════════════════════

\begin{frame}{Ablauf}

\begin{columns}
\begin{column}{0.5\textwidth}
\textbf{So funktioniert es:}
\begin{enumerate}
    \item QR-Code scannen
    \item Lernpfad starten
    \item Punkte sammeln
    \item Badges verdienen
\end{enumerate}
\end{column}

\begin{column}{0.5\textwidth}
\begin{block}{Ziel}
\centering
\textbf{80+ Punkte} erreichen!
\end{block}

\vspace{3mm}

\textbf{5 Badges möglich:}
\begin{itemize}
    \item Vocabulario-Meister
    \item Gramática-Profi
    \item Perfecto-Experte
    \item Pronomen-Held
    \item Präpositionen-Profi
\end{itemize}
\end{column}
\end{columns}

\end{frame}

% ═══════════════════════════════════════════════════════════════════════════
% QR-CODE
% ═══════════════════════════════════════════════════════════════════════════

\begin{frame}{Lernpfad starten}

\begin{center}

\vspace{5mm}

\qrcode[height=50mm]{https://devbox42.github.io/sciencesim/spanisch/kl09/unidad-02-mi-nueva-casa/DS-04/LP-repaso-U1-U2.html}

\vspace{8mm}

{\large\color{gray}Scanne den QR-Code mit deinem Gerät}

\end{center}

\end{frame}

% ═══════════════════════════════════════════════════════════════════════════
% HILFE
% ═══════════════════════════════════════════════════════════════════════════

\begin{frame}{Hilfe im Lernpfad}

\begin{columns}
\begin{column}{0.6\textwidth}
\textbf{Wenn du nicht weiterkommst:}
\begin{itemize}
    \item Klicke auf den \textbf{Hilfe-Button}
    \item Du bekommst die wichtigste Regel angezeigt
    \item Bei Fehlern gibt es \textbf{Erklärungen}
\end{itemize}

\vspace{5mm}

\textbf{Checklist nutzen:}
\begin{itemize}
    \item Schau auf deine Checklist
    \item Dort stehen alle Regeln
\end{itemize}
\end{column}

\begin{column}{0.4\textwidth}
\begin{block}{Tipp}
Lies die Regel-Boxen im Lernpfad aufmerksam durch, bevor du die Übung machst!
\end{block}
\end{column}
\end{columns}

\end{frame}

% ═══════════════════════════════════════════════════════════════════════════
% ARBEITSAUFTRAG
% ═══════════════════════════════════════════════════════════════════════════

\begin{frame}{Arbeitsauftrag}

\begin{block}{Jetzt: Lernpfad bearbeiten (80 Min)}
\begin{enumerate}
    \item Öffne den Lernpfad (QR-Code)
    \item Arbeite alle Übungen durch
    \item Sammle so viele Punkte wie möglich
    \item Am Ende: Ergebnis merken!
\end{enumerate}
\end{block}

\vspace{5mm}

\begin{columns}
\begin{column}{0.5\textwidth}
\textbf{Bei technischen Problemen:}\\
Melde dich -- du bekommst das Arbeitsblatt
\end{column}

\begin{column}{0.5\textwidth}
\textbf{Fertig vor der Zeit?}\\
Wiederhole schwache Bereiche oder hilf anderen
\end{column}
\end{columns}

\end{frame}

% ═══════════════════════════════════════════════════════════════════════════
% ABSCHLUSS
% ═══════════════════════════════════════════════════════════════════════════

\begin{frame}{Wie lief es?}

\begin{center}

\textbf{Wer hat 80+ Punkte erreicht?}

\vspace{10mm}

\textbf{Welche Badges habt ihr bekommen?}

\vspace{10mm}

\textbf{Was war schwierig?}

\end{center}

\end{frame}

% ═══════════════════════════════════════════════════════════════════════════
% AUSBLICK
% ═══════════════════════════════════════════════════════════════════════════

\begin{frame}{Vor der Prüfung}

\begin{block}{Noch zu tun:}
\begin{itemize}
    \item Checklist durchgehen -- alles abgehakt?
    \item Schwache Bereiche nochmal üben
    \item Unregelmäßige Partizipien lernen
\end{itemize}
\end{block}

\vspace{5mm}

\begin{center}
{\Large\color{spanisch}\textbf{¡Buena suerte en el examen!}}
\end{center}

\end{frame}

\end{document}
